\subsection{Further results on compact metrizable topological spaces}

We endow the two-element set \(\{0,1\}\) with the discrete topology and the set \(\{0,1\}^{N}\) with the product topology. The topological space \(\{0,1\}^{N}\) is compact (TG, I, p. 63, th. 3), metrizable (TG, IX, p. 15, cor. 2), nonempty, totally disconnected (TG, I, p. 84, prop. 10), and has no isolated point.

PROPOSITION 14. --- Every compact, metrizable, nonempty, totally disconnected topological space with no isolated point is homeomorphic to \(\{0,1\}^{N}\) .

Let X be such a topological space. Endow it with a distance compatible with its topology. Since the space X is compact, it is complete for this distance (TG, II, p. 27, th. 1).

Lemma 3. --- Let \(\varepsilon\) be a real number \(>0\) . There exists an integer \(m \geqslant 1\) such that, for every integer \(n \geqslant m\) , X admits a partition consisting of \(n\) nonempty open and closed sets of diameter \(\leqslant \varepsilon\) .

Every point of X admits an open and closed neighborhood of diameter \(\leqslant\varepsilon\) (TG, II, p. 32, corollary of prop. 6). Since the space X is compact, it possesses a finite covering by such sets. Choose one of these, \((\mathrm{U}_{i})_{1\leqslant i\leqslant m}\) , for which m is minimal. We have \(m\geqslant1\) since X is not empty. For \(1\leqslant i\leqslant m\) , denote by \(V_{i}\) the intersection of \(U_{i}\) and the

\(X\setminus U_k\) for \(k < i\). Then \((V_i)_{1 \leqslant i \leqslant m}\) is a partition of X, consisting of m nonempty open and closed sets of diameter \(\leqslant \varepsilon\).

Since X has no isolated point, every nonempty open and closed subset V of X contains at least two points. Moreover, since X is compact and totally disconnected, V is the union of two disjoint nonempty open and closed subsets (loc. cit.). It follows, by induction, that for every integer \(n \geqslant m\) , X admits a partition consisting of n nonempty open and closed sets of diameter \(\leqslant \varepsilon\) .

Let us now finish the proof of Prop. 14. Every nonempty open and closed subspace of X is a compact, totally disconnected metric space with no isolated point. Lemma 3 therefore allows us to construct by induction a sequence \((\mathrm{J}_{n})_{n\geqslant0}\) of finite sets and, for every \(n\geqslant0\), a map \(\varphi_{n}\) from the set \(C_{n}=J_{0}\times\cdots\times J_{n}\) into the set of nonempty open and closed subsets of X of diameter \(\leqslant2^{-n}\), such that:

\begin{enumerate}
\def\labelenumi{(\roman{enumi})}
\item
  For every \(n \geqslant 0\) , there exists an integer \(m_n \geqslant 1\) such that \(\mathrm{Card}(\mathrm{J}_n) = 2^{m_n}\) ;
\item
  The family \((\varphi_0(c))_{c\in \mathrm{C}_0}\) is a partition of \(X\) ;
\item
  For every \(n \geqslant 0\) and every \(c \in \mathrm{C}_n\) , the family \((\varphi_{n+1}(c,j))_{j \in \mathrm{J}_{n+1}}\) is a partition of \(\varphi_n(c)\) .
\end{enumerate}

Denote by \(p_n\) the canonical projection of \(C_{n+1}\) onto \(C_n\). The sequence \(C = (C_n, p_n)_{n \geqslant 0}\) is a sieve (TG, IX, p. 63, def. 8). The topological space associated with this sieve (TG, IX, p. 63) is identified with the topological space J, the product of the discrete topological spaces \(J_n\). The sieve C and the sequence of maps \((\varphi_n)_{n \geqslant 0}\) define a strict sieving of the metric space X (TG, IX, p. 63 and p. 64). The map \(f: J \to X\) deduced from this sieving is continuous and bijective (TG, IX, p. 65). Since the topological space J is compact (TG, I, p. 63, th. 3) and X is Hausdorff, \(f\) is a homeomorphism (TG, I, p. 63, cor. 2). Since \(J_n\) is homeomorphic to \(\{0, 1\}^{m_n}\) for all \(n \geqslant 0\), J is homeomorphic to \(\{0, 1\}^N\) (TG, I, p. 25, prop. 2).

Example. --- Let K be the Cantor ternary set (TG, IV, p. 9, example). For every \(n \geqslant 0\) , set \(J_{n} = \{0,1\}\) and define a map \(K_{n}\) from the set \(C_{n} = J_{0} \times \cdots \times J_{n}\) into the set of closed intervals of {[}0,1{]} as follows: we set \(K_{0}(0) = [0, \frac{1}{3}]\) and \(K_{0}(1) = [\frac{2}{3}, 1]\) ; for every \(n \geqslant 0\) and every \(c \in C_{n}\) , \(K_{n+1}(c,0)\) and

\(K_{n+1}(c,1)\) are respectively the "left third" and the "right third" of \(K_n(c)\). If \(c = (j_0,j_1,\ldots,j_n) \in C_n\), \(K_n(c)\) is the interval denoted \(K_{n,p}\) in loc. cit., with \(p = 2^n j_0 + 2^{n-1}j_1 + \cdots + j_n + 1\), it is also the interval \([a,a + \frac{1}{3^{n+1}}]\), where \(a = 2(\frac{j_0}{3} + \frac{j_1}{3^2} + \cdots + \frac{j_n}{3^{n+1}})\). For \(n \geqslant 0\) and \(c \in C_n\), set \(\varphi_n(c) = K_n(c) \cap K\). The family \((\varphi_n(c))_{c \in C_n}\) is a partition of K formed by closed sets. These sets are therefore also open in K; they are nonempty and of diameter \(\frac{1}{3^{n+1}}\), because the endpoints of the intervals \(K_n(c)\) belong to K. The sequence \(C = (C_n, p_n)_{n \geqslant 0}\), where \(p_n: C_{n+1} \to C_n\) is the canonical projection \((c,j) \mapsto c\), is a sieve, and the topological space associated with this sieve is identified with \(\{0,1\}^N\). The sieve C and the sequence of maps \((\varphi_n)_{n \geqslant 0}\) define a strict sieving of the metric space K. The map \(f: \{0,1\}^N \to K\) deduced from this sieving is a homeomorphism, given by the formula

\[
f ((j _ {n}) _ {n \geqslant 0}) = 2 \sum_ {n = 0} ^ {\infty} \frac {j _ {n}}{3 ^ {n + 1}}.
\]

COROLLARY. --- Let X be a metrizable, compact, and nonempty topological space. There exists a continuous and surjective map from \(\{0,1\}^{N}\) to X.

Every compact and metrizable topological space is homeomorphic to a subspace, necessarily closed, of the topological space \(I^{N}\) (TG, IX, p. 18, prop. 12 and p. 21, prop. 16). Let A be a closed subspace of \(I^{N}\) and let h be a homeomorphism from A onto X.

Let us set \(K = \{0,1\}^{N}\) and, for \(\alpha = (a_n) \in K\), set \(f(\alpha) = \sum_{n=0}^{\infty} a_n 2^{-n-1}\); we have \(f(\alpha) \in I\). The map \(f: K \to I\) thus defined is surjective (TG, IV, p. 42) and continuous. Indeed, if two elements \(\alpha\) and \(\beta\) of K have the same coordinates of index \(< n\), we have \(|f(\beta) - f(\alpha)| \leqslant 2^{-n}\). Denote by \(g\) the map \((\alpha_n) \mapsto (f(\alpha_n))\) from \(K^N\) into \(I^N\); it is continuous and surjective. The topological space \(K^N\) is compact (TG, I, p. 63, th. 3), metrizable (TG, IX, p. 15, cor. 2), and totally disconnected (TG, I, p. 84, prop. 10), and the same holds for its closed subspace \(\bar{g}^{1}(A)\); the latter is nonempty since the maps \(g\) and \(h\) are surjective.

Then, the space \(\bar{g}^{1}(A) \times K\) is homeomorphic to \(\{0,1\}^{N}\) (prop. 14), since it is a compact, metrizable, totally disconnected topological space without isolated points, and the map \((x,y) \mapsto h(g(x))\) of \(\bar{g}^{1}(A) \times K\) into X is continuous and surjective.

